\lecture{9}{23. Marts 2026}{Thermal strains}

\subsection{Thermal Strains}
Temperature changes may be a source of stresses in elastic systems.
It is reasonable to assume that a linear relationship exists between the temperature difference from the datum value and the corresponding strain
\begin{equation*}
  \epsilon_{(ii)} (T) = \alpha \left( T - T_0 \right)
\end{equation*}
in which $\epsilon_{(ii)} (T)$ is the linear strain due to the temperature difference $T - T_0$, $T_0$ is the datum temperature and $\alpha$ is the coefficient of thermal expansion.

Using this, we can write
\begin{equation*}
  \epsilon_{ij} = \frac{1}{E} \left( \left( 1 + \nu \right)\sigma_{ij} - \nu \delta_{ij} \sigma_{kk} \right) + \alpha \delta_{ij} \left( T - T_0 \right)
\end{equation*}
and
\begin{equation*}
  \sigma_{ij} = 2\mu\epsilon_{ij} + \lambda \delta_{ij} \epsilon_{kk} - \alpha \delta_{ij} \left( 3 \lambda + 2\mu \right) \left( T - T_0 \right)
.\end{equation*}
